% Week4: SNN Training 讲稿
% 对应 PPT: Week4_SNN Training.pptx
\documentclass[journal,12pt,onecolumn]{IEEEtran}
\IEEEoverridecommandlockouts
\usepackage{cite}
\usepackage{amsmath, amssymb, bm, color, graphicx, enumitem}
\usepackage{amsmath,amssymb,amsfonts}
\usepackage{algorithmic}
\usepackage{algorithm}
\usepackage{graphicx}
\usepackage{booktabs}
\usepackage{textcomp}
\usepackage{xcolor}
\usepackage{subfigure}
\graphicspath{{./img/}}
\usepackage{tabularx}
\usepackage{makecell}
\usepackage{caption2}
\usepackage{color}

\usepackage[UTF8]{ctex}

\definecolor{myred}{RGB}{255, 67, 20}
\definecolor{myblue}{RGB}{51, 51, 255}
\definecolor{mygreen}{RGB}{0, 128, 0}

\begin{document}
\title{Week4: SNN 在线训练方法\\OTTT 与 NDOT 讲稿}
\author{Tao}
\date{\today}
\maketitle

\begin{abstract}
本文为 Week4 ``SNN Training'' 课程的讲稿内容，对应于 Week4\_SNN Training.pptx。在前三周已讲授 SNN 前向动力学、BPTT 与代理梯度的基础上，本周系统介绍从 BPTT 到在线训练的方法论跃迁，重点讲述 OTTT（Online Training Through Time）和 NDOT（Neuronal Dynamics-based Online Training）两种代表性在线训练算法，包括其数学推导、核心近似、递推实现及对比分析。
\end{abstract}

% ============================================================
% Slide 1: 课程回顾与本周概述
% ============================================================
\section{课程回顾与本周概述}

\subsection*{前情回顾}

在前三周的课程中，我们已经讲授了以下内容：

\begin{itemize}
    \item \textbf{Week 1}：SNN 的基本概念——LIF 神经元模型、脉冲生成机制、与 ANN 的核心差异。
    \item \textbf{Week 2}：BPTT（Backpropagation Through Time）的完整推导——将 SNN 沿时间展开为深度计算图，利用链式法则计算梯度。核心结论：BPTT 梯度精确，但需要存储全部时间步的中间状态（O($T$) 内存），且无法在线更新。
    \item \textbf{Week 3}：STBP（Spatio-Temporal Backpropagation）与代理梯度（Surrogate Gradient）——用光滑函数（如矩形函数、sigmoid 函数）的导数替代不可导的 Heaviside 阶跃函数导数，使梯度能够传播。代理梯度的引入使得 BPTT 在 SNN 中实际可用，但未能解决内存和在线更新的瓶颈。
\end{itemize}

\subsection*{本周内容}

\begin{itemize}
    \item \textbf{核心问题}：如何在保持梯度精度的前提下，实现 O(1) 内存的在线训练？
    \item \textbf{方法论主线}：
    \[
    \text{BPTT（离线，O($T$)）} \;\longrightarrow\; \text{OTTT（在线，O(1)，近似）} \;\longrightarrow\; \text{NDOT（在线，O(1)，物理）}
    \]
    \item \textbf{关键思想}：将 BPTT 中需要时间回溯的连乘项 $\prod \epsilon[i]$，简化为可前向递推的衰减累积，从而打破 O($T$) 内存限制。
\end{itemize}

\paragraph{讲稿：}
今天我们进入第四周的课程——SNN 的在线训练方法。前三周我们打下了扎实的基础：从 LIF 神经元、到 BPTT 的完整推导、再到代理梯度。今天我们要解决的核心问题是：BPTT 虽然精确，但内存随序列长度线性增长，这对边缘设备和神经形态芯片来说是不可接受的。那么，有没有办法让 SNN 逐个时间步地在线学习、内存不随时间增长呢？答案是有的——OTTT 和 NDOT 就是两种代表性的解决方案。

% ============================================================
% Slide 2: BPTT 的核心瓶颈回顾
% ============================================================
\section{BPTT 的核心瓶颈回顾}

\subsection*{LIF 神经元离散动力学}

在进入在线方法之前，我们先回顾 BPTT 的核心公式。LIF 神经元的离散动力学为：

\begin{equation}
\mathbf{u}^{l}[t+1] = \lambda (\mathbf{u}^{l}[t] - V_{th}\mathbf{s}^{l}[t]) + \mathbf{W}^{l}\mathbf{s}^{l-1}[t+1]
\label{eq.lif}
\end{equation}

其中：
\begin{itemize}
    \item $\mathbf{u}^{l}[t]$：第 $l$ 层 $t$ 时刻的膜电位
    \item $\mathbf{s}^{l}[t] = H(\mathbf{u}^{l}[t] - V_{th})$：脉冲输出（Heaviside 阶跃函数）
    \item $\lambda$：泄漏常数（$0 < \lambda \leq 1$）
    \item $V_{th}$：发放阈值
    \item $\mathbf{W}^{l}$：待训练的权重矩阵
\end{itemize}

\subsection*{BPTT 完整梯度公式}

BPTT 对权重 $\mathbf{W}^l$ 的梯度为：

\begin{equation}
\boxed{
\begin{aligned}
\frac{\partial L}{\partial \mathbf{W}^l}
&=
\sum_{t=1}^T
\frac{\partial L}{\partial \mathbf{s}^{l+1}[t]}
{\color{red}{\frac{\partial \mathbf{s}^{l+1}[t]}{\partial \mathbf{u}^{l+1}[t]}}}
\Bigg(
\frac{\partial \mathbf{u}^{l+1}[t]}{\partial \mathbf{W}^l}
\\
&\quad +
\sum_{\tau < t}
\prod_{i=t-1}^{\tau}
\Big(
{\color[rgb]{0,0.69,0.31}{\frac{\partial \mathbf{u}^{l+1}[i+1]}{\partial \mathbf{u}^{l+1}[i]}}}
+
{\color{cyan}{\frac{\partial \mathbf{u}^{l+1}[i+1]}{\partial \mathbf{s}^{l+1}[i]}}}
{\color{red}{\frac{\partial \mathbf{s}^{l+1}[i]}{\partial \mathbf{u}^{l+1}[i]}}}
\Big)
\frac{\partial \mathbf{u}^{l+1}[\tau]}{\partial \mathbf{W}^l}
\Bigg)
\end{aligned}
}
\label{eq.bptt}
\end{equation}

\subsection*{BPTT 的两大瓶颈}

\begin{enumerate}
    \item \textbf{内存瓶颈}：连乘项 $\prod_{i=t-1}^{\tau}$ 需要存储全部 $T$ 个时刻的 $\mathbf{u}$ 和 $\mathbf{s}$，内存复杂度 O($T$)。
    \item \textbf{离线瓶颈}：必须等待整个长度为 $T$ 的序列处理完毕才能计算梯度并更新权重，无法实现在线（流式）学习。
\end{enumerate}

这两大瓶颈的根本原因都指向同一个结构——公式~\eqref{eq.bptt} 中橙色的连乘项，它要求我们沿着时间轴反向追溯，访问所有历史状态。

\paragraph{讲稿：}
我们先来回顾 BPTT 的梯度公式。注意看这个公式的结构：外层是对时间 $t$ 的求和，这是因为权重在所有时间步共享。括号里有两项——当前直连项和历史回溯项。当前直连项很简单，等于当前输入脉冲 $\mathbf{s}^{l}[t]$。问题出在历史回溯项，它包含一个从 $\tau$ 到 $t-1$ 的连乘。这个连乘里有两部分：绿色的是泄漏路径 $\lambda$，青色×红色的是重置路径。因为连乘涉及整个历史区间，反向传播时必须把所有中间状态存在内存里，这就导致了 O($T$) 的内存复杂度。今天我们要讲的所有方法，本质上都是在"如何打破这个连乘依赖"上做文章。

% ============================================================
% Slide 3: 从 BPTT 到在线学习的核心思路
% ============================================================
\section{从 BPTT 到在线学习：核心思路}

\subsection*{为什么连乘是瓶颈？}

BPTT 梯度中的连乘结构：
\begin{equation}
\prod_{i=t-1}^{\tau} \epsilon[i], \quad \epsilon[i] = {\color[rgb]{0,0.69,0.31}{\lambda}} + {\color{cyan}{(-\lambda V_{th})}}{\color{red}{\frac{\partial s}{\partial u}}}
\label{eq.epsilon}
\end{equation}

这个结构在反向传播时天然需要访问所有中间状态——这是一个因果链：要计算 $t$ 时刻的梯度，必须通过每一个中间时间步反向传递误差信号。

\subsection*{核心突破点}

关键在于观察：
\[
\frac{\partial \mathbf{u}^{l+1}[\tau]}{\partial \mathbf{W}^l} = \mathbf{s}^{l}[\tau]
\]

即膜电位对权重的直接偏导恰好等于输入脉冲。如果时序核能够退化为一个纯衰减因子，那么历史贡献就变成了：

\[
\sum_{\tau \leq t} \lambda^{t-\tau} \mathbf{s}^{l}[\tau]
\]

这个式子可以通过前向递推高效计算：
\[
\hat{\mathbf{a}}^{l}[t] = \lambda \hat{\mathbf{a}}^{l}[t-1] + \mathbf{s}^{l}[t]
\]

只需要维护一个大小为 $n$ 的向量 $\hat{\mathbf{a}}$，内存从 O($T$) 骤降至 O(1)！

\paragraph{讲稿：}
大家可能已经注意到了，BPTT 的瓶颈完全集中在那个时间连乘上。如果我们能把这个连乘简化，是不是就能实现在线学习呢？这个思路就是 OTTT 的出发点。我们注意到一个关键的恒等式：膜电位对权重的偏导正好是输入脉冲。如果时序核能退化成简单的指数衰减 $\lambda^{t-\tau}$，那么整个历史回溯项就变成了一个前向可递推的量 $\hat{a}[t]$。这就把 O($T$) 的内存需求削减到了 O(1)。关键问题是：怎么处理那个复杂的 $\epsilon[i]$，让它退化成 $\lambda$？

% ============================================================
% Slide 4: OTTT 的核心近似——丢弃重置路径
% ============================================================
\section{OTTT：核心近似与推导}

\subsection*{OTTT 的核心近似}

OTTT 的核心理念非常简洁：在时序连乘中，不对红色项施加代理梯度，而是直接取其真实导数（几乎处处为 0）：

\begin{equation}
{\color{cyan}{\frac{\partial \mathbf{u}[i+1]}{\partial \mathbf{s}[i]}}}
{\color{red}{\frac{\partial \mathbf{s}[i]}{\partial \mathbf{u}[i]}}}
\approx 0
\quad\Longrightarrow\quad
\prod_{i=\tau}^{t-1}
\left({\color[rgb]{0,0.69,0.31}{\lambda \mathbf{I}}} + 0\right)
= \lambda^{t-\tau} \mathbf{I}
\label{eq.ottt_approx}
\end{equation}

\subsection*{简化后的梯度公式}

经过 OTTT 近似后，BPTT 的梯度退化为：

\begin{equation}
\frac{\partial L}{\partial \mathbf{W}^l}
=
\sum_{t=1}^T
\frac{\partial L}{\partial \mathbf{s}^{l+1}[t]}
{\color{red}{\frac{\partial \mathbf{s}^{l+1}[t]}{\partial \mathbf{u}^{l+1}[t]}}}
\left(
\sum_{\tau \leq t} \lambda^{t-\tau}
\frac{\partial \mathbf{u}^{l+1}[\tau]}{\partial \mathbf{W}^l}
\right)
\label{eq.ottt_gradient}
\end{equation}

注意，红色项仅在空间传播入口处保留（仍需代理梯度），但在时序回溯中已被完全消除。

\subsection*{定义追踪变量}

定义突触前活动踪迹：
\begin{equation}
\boxed{
\hat{\mathbf{a}}^{l}[t] \triangleq \sum_{\tau=1}^{t} \lambda^{t-\tau} \mathbf{s}^{l}[\tau]
}
\label{eq.trace}
\end{equation}

利用 $\frac{\partial \mathbf{u}^{l+1}[\tau]}{\partial \mathbf{W}^l} = \mathbf{s}^{l}[\tau]$，梯度可重写为：
\begin{equation}
\nabla_{\mathbf{W}^l}L = \sum_{t=1}^T \mathbf{g}_{\mathbf{u}^{l+1}}[t] \left(\hat{\mathbf{a}}^{l}[t]\right)^\top
\label{eq.ottt_final}
\end{equation}
其中 $\mathbf{g}_{\mathbf{u}^{l+1}}[t] = \left( \frac{\partial L}{\partial \mathbf{s}^{l+1}[t]} {\color{red}{\frac{\partial \mathbf{s}^{l+1}[t]}{\partial \mathbf{u}^{l+1}[t]}}} \right)^\top$ 为空间梯度。

\subsection*{前向递推（在线计算的关键）}

$\hat{\mathbf{a}}^{l}[t]$ 可以通过前向递推在线计算：
\begin{equation}
\boxed{
\hat{\mathbf{a}}^{l}[t] = \lambda \hat{\mathbf{a}}^{l}[t-1] + \mathbf{s}^{l}[t]
}
\label{eq.ottt_recurrence}
\end{equation}

这一递推仅依赖上一时刻的 $\hat{\mathbf{a}}$ 和当前时刻的脉冲，完全不需要存储或访问更早期的历史状态。

\paragraph{讲稿：}
OTTT 的核心近似非常大胆——它直接说"重置路径那一项（青色×红色）对时序传播的贡献可以忽略不计"。为什么可以这样近似？因为红色项 $\partial s / \partial u$ 的真实导数几乎处处为零，只在阈值附近有一个无穷大的尖峰。既然我们本来就用代理梯度去近似它，那在时间传播中干脆把它置零，只保留确定性的泄漏项 $\lambda$，这不是更一致吗？

这个近似带来了巨大的实际好处。连乘退化为 $\lambda^{t-\tau}$，是一个固定的指数衰减，与神经元的具体发放状态无关。然后我们定义 $\hat{a}[t]$ 为历史脉冲的指数加权和——这恰好就是突触前活动的"资格迹"。最关键的是，$\hat{a}[t]$ 可以用递推式 $\hat{a}[t] = \lambda \hat{a}[t-1] + s[t]$ 前向计算，只需要存一个向量，内存 O(1)，完美实现在线学习。

最终梯度变成了空间梯度与追踪变量的外积 $\mathbf{g} \cdot \hat{a}^\top$，形式和 BPTT 一致，但历史信息全部被 $\hat{a}$ 捕获，不再需要反向时间传播。

% ============================================================
% Slide 5: OTTT 的简化效果与局限
% ============================================================
\section{OTTT 的简化效果与局限分析}

\subsection*{OTTT 的简化效果}

\begin{table}[h]
\centering
\caption{BPTT vs OTTT 对比}
\begin{tabular}{lcc}
\toprule
\textbf{特性} & \textbf{BPTT} & \textbf{OTTT} \\
\midrule
内存复杂度 & O($T$) & O(1) \\
在线更新     & 否   & 是 \\
时间传播核   & $\prod(\lambda + \text{reset})$ & $\lambda^{t-\tau}$ \\
代理梯度使用 & 时间+空间 & 仅空间 \\
追踪变量     & 无   & $\hat{a}[t] = \lambda\hat{a}[t-1] + s[t]$ \\
\bottomrule
\end{tabular}
\label{tab.bptt_ottt}
\end{table}

\subsection*{OTTT 的局限性}

尽管 OTTT 在工程上取得了巨大成功（首次实现了 ImageNet 级别的 SNN 在线训练），但它的近似过于粗暴：

\begin{enumerate}
    \item \textbf{固定衰减 $\lambda$}：所有神经元在任何时刻都使用相同的衰减常数，但真实神经元中，膜电位的衰减受充电、泄漏、发放、重置的共同影响，是一个动态过程。
    \item \textbf{忽略重置路径}：当神经元发放 spike 后，膜电位会被重置（$u \leftarrow u - V_{th}$），这一突变会影响后续所有时间步的膜电位动态。OTTT 完全忽略了这个影响。
    \item \textbf{精度损失}：在需要精确时序信息的任务上（如 DVS 事件流处理、时序分类），固定衰减的近似可能导致精度下降。
\end{enumerate}

\paragraph{讲稿：}
OTTT 实现了从 O($T$) 到 O(1) 的飞跃，这是一个很大的工程贡献。但我们需要诚实地面对它的局限：所有神经元用同一个固定的 $\lambda$，这意味着无论神经元当前处于什么状态——正在充电、刚刚放电、还是处于静息——它的"历史衰减"都是一样的。这显然不符合真实的神经元动力学。那么，有没有办法用一个动态的、依赖神经元状态的衰减系数来替代固定的 $\lambda$ 呢？这就是 NDOT 要做的事情。

% ============================================================
% Slide 6: NDOT 的物理重参数化
% ============================================================
\section{NDOT：基于神经元动力学的重参数化}

\subsection*{从连续 LIF 方程出发}

NDOT 的核心思想是不再删除任何项，而是利用神经元自身的物理方程来推导时序核。我们从连续时间的 LIF 微分方程开始：

\begin{equation}
u'(t) = -u(t) + I(t)
\label{eq.cont_lif}
\end{equation}

其中 $I(t)$ 为输入电流。这个方程描述了膜电位在连续时间下的演化——泄漏（$-u(t)$）和充电（$I(t)$）。

\subsection*{定义动态时序依赖 $e(t)$}

NDOT 定义了一个核心量——相邻时刻膜电位的变化率之比：

\begin{equation}
e(t) \triangleq \frac{\partial u(t+1)}{\partial u(t)}
= \frac{u'(t+1)}{u'(t)}
\label{eq.e_def}
\end{equation}

这一步的直觉是：如果我们能知道膜电位在两个相邻时刻之间的变化率之比，就能知道前一时刻的状态对后一时刻的影响有多大。

\subsection*{消去输入电流 $I(t)$}

将 LIF 方程代入：
\begin{equation}
e(t) = \frac{-u(t+1) + I(t+1)}{-u(t) + I(t)}
\label{eq.e_substitute}
\end{equation}

\subsection*{离散化并化简}

在离散 LIF 动力学中，输入电流满足：
\[
I[t+1] = u[t+1] - \lambda(u[t] - V_{th}s[t])
\]

代入并化简，得到 NDOT 的核心闭式解：

\begin{equation}
\boxed{
\mathbf{e}^{l}[t] =
\frac{\mathbf{u}^{l}[t] - V_{th}\mathbf{s}^{l}[t]}
     {\mathbf{u}^{l}[t-1] - V_{th}\mathbf{s}^{l}[t-1]}
}
\label{eq.ndot_e}
\end{equation}

这个表达式的物理含义非常清晰：它是两个相邻时刻"重置后膜电位"的比值。分子和分母分别是当前时刻和前一时刻发放重置后的残余膜电位。注意到，它完全不包含代理梯度——纯粹由前向动力学状态决定。

\paragraph{讲稿：}
NDOT 的思路和 OTTT 完全不同。OTTT 是"我删掉一项"，NDOT 是"我用物理来重新推导这一项"。NDOT 从连续时间的 LIF 微分方程出发：$u'(t) = -u(t) + I(t)$。它没有用 $\lambda$ 作为固定的离散衰减，而是定义 $e(t) = u'(t+1) / u'(t)$——即膜电位变化率的比值——来衡量相邻时刻的依赖性。

关键的推导技巧是利用隐函数定理：在连续系统中，$u(t+1)$ 对 $u(t)$ 的偏导等于两者对时间导数之比。然后把 LIF 方程代入，消去输入电流 $I(t)$，再离散化代入离散 LIF 方程——最终得到了一个极其漂亮的闭式解：$\mathbf{e}[t]$ 等于相邻时刻重置后膜电位的比值。

这个 $e[t]$ 的值会在每个时间步、对每个神经元独立计算。如果某个神经元刚刚发放，它的 $e[t]$ 会比 $\lambda$ 小；如果神经元长时间没有发放，$e[t]$ 接近 $\lambda$。这就捕捉了 OTTT 完全丢失的信息——神经元状态对时序传播的动态影响。

% ============================================================
% Slide 7: NDOT 的动态时序核与望远镜化简
% ============================================================
\section{NDOT 的动态时序核与望远镜化简}

\subsection*{定义动态时序传播核}

与 OTTT 的 $\lambda^{t-\tau}$ 对应，NDOT 定义了动态时序传播核：

\begin{equation}
\mathbf{P}_{k,t}^{l} \triangleq \prod_{i=k}^{t-1} \mathbf{e}^{l}[i]
= \mathbf{e}^{l}[k] \odot \mathbf{e}^{l}[k+1] \odot \cdots \odot \mathbf{e}^{l}[t-1]
\label{eq.P_def}
\end{equation}

其中 $\odot$ 为逐元素乘法（Hadamard product）。

\subsection*{望远镜化简（Telescoping Cancellation）}

这是 NDOT 最优雅的数学性质之一。将 $\mathbf{e}[i]$ 的分式形式代入 $\mathbf{P}_{k,t}$：

\begin{equation}
\begin{aligned}
\mathbf{P}_{k,t}^{l}
&=
\frac{u[k] - V_{th}s[k]}{u[k-1] - V_{th}s[k-1]}
\odot
\frac{u[k+1] - V_{th}s[k+1]}{u[k] - V_{th}s[k]}
\odot \cdots \odot
\frac{u[t-1] - V_{th}s[t-1]}{u[t-2] - V_{th}s[t-2]}
\\
&=
\frac{\mathbf{u}^{l}[t-1] - V_{th}\mathbf{s}^{l}[t-1]}
     {\mathbf{u}^{l}[k-1] - V_{th}\mathbf{s}^{l}[k-1]}
\end{aligned}
\label{eq.telescoping}
\end{equation}

中间所有相邻项相互约去，最终只剩下首尾两项的比值！这意味着，虽然 $\mathbf{P}_{k,t}$ 在定义上是一个连乘，实际计算时可以一步到位，无需逐项相乘。

\subsection*{追踪变量的递推}

利用 $\mathbf{P}_{k,t}$ 的递推性质，NDOT 同样可以定义前向递推的追踪变量：

\begin{equation}
\boxed{
\hat{\mathbf{a}}^{l}[t] = \mathbf{e}^{l}[t-1] \odot \hat{\mathbf{a}}^{l}[t-1] + \mathbf{s}^{l}[t]
}
\label{eq.ndot_recurrence}
\end{equation}

\subsection*{最终梯度公式}

与 OTTT 形式完全一致：
\begin{equation}
\boxed{
\nabla_{\mathbf{W}^l} \mathcal{L}
=
\sum_{t=1}^{T}
\mathbf{g}_{\mathbf{u}^l}[t]
\left(
\hat{\mathbf{a}}^{l-1}[t]
\right)^{\top}
}
\label{eq.ndot_final}
\end{equation}

\paragraph{讲稿：}
NDOT 的 $\mathbf{P}_{k,t}$ 是对 OTTT 的 $\lambda^{t-k}$ 的直接推广——从标量指数衰减推广为逐元素的动态比值连乘。但这里有一个非常漂亮的数学性质：当每个步骤都是分式比值时，连乘会产生望远镜效应——中间所有相邻的分子分母相互约去，最终 $\mathbf{P}_{k,t}$ 仅等于首尾两项的比值。这意味着我们不实际上需要计算那个连乘——计算效率和 OTTT 一样高效。

更进一步，追踪变量的递推形式与 OTTT 完全一致：$\hat{a}[t] = e[t-1] \odot \hat{a}[t-1] + s[t]$，唯一的区别就是把固定的 $\lambda$ 换成了动态的 $e[t-1]$。最终梯度也是 $\mathbf{g} \cdot \hat{a}^\top$ 的形式。所以 NDOT 在保持和 OTTT 完全相同的 O(1) 内存和计算复杂度的前提下，用更精确的物理量替换了近似常数。

% ============================================================
% Slide 8: OTTT vs NDOT 核心对比
% ============================================================
\section{OTTT vs NDOT 核心对比}

\begin{table}[h]
\centering
\caption{OTTT 与 NDOT 核心对比}
\begin{tabular}{lcc}
\toprule
\textbf{维度} & \textbf{OTTT} & \textbf{NDOT} \\
\midrule
时序核（单步）  & $\epsilon[i] = \lambda$ & $\mathbf{e}[i] = \frac{u[i]-V_{th}s[i]}{u[i-1]-V_{th}s[i-1]}$ \\
历史传播核    & $\lambda^{t-k}$ & $\mathbf{P}_{k,t} = \prod \mathbf{e}[i]$（望远镜化简） \\
追踪变量递推  & $\hat{a}[t] = \lambda\hat{a}[t-1] + s[t]$ & $\hat{a}[t] = \mathbf{e}[t-1]\odot\hat{a}[t-1] + s[t]$ \\
额外存储      & 无 & 需存 $u[t]-V_{th}s[t]$ \\
方法论本质    & 近似删除 & 物理重参数化 \\
内存复杂度    & O(1) & O(1) \\
\bottomrule
\end{tabular}
\label{tab.ottt_ndot}
\end{table}

\subsection*{一句话总结}

\begin{quote}
OTTT 是把 BPTT 中的重置路径\textbf{删掉}，用固定 $\lambda$ 近似时序依赖；NDOT 是用连续的神经元物理方程\textbf{重新推导}时序依赖，用动态比值 $\mathbf{e}[t]$ 替代固定 $\lambda$，在不增加计算成本的前提下获得更精确的梯度估计。
\end{quote}

\paragraph{讲稿：}
现在是总结对比的时候。OTTT 和 NDOT 的外在形式惊人地相似——相同的递推结构、相同的 O(1) 内存、相同的 $\mathbf{g}\cdot\hat{a}^\top$ 最终形式。但它们的本质差异在于方法论：OTTT 走的是"近似简化"的路线——我把一项设为零，问题就简单了；NDOT 走的是"物理精确化"的路线——我不删除任何东西，而是从连续动力学重新推导出正确的依赖关系。

结果就是，NDOT 在相同的计算成本下获得了更高的精度。在 CIFAR-100 上，NDOT 比 OTTT 提升了大约 5.5 个百分点。这是一个相当大的提升，而且"免费"——没有增加额外的计算负担。

% ============================================================
% Slide 9: 实验结果对比
% ============================================================
\section{实验结果对比}

\subsection*{OTTT vs NDOT 精度对比}

\begin{table}[h]
\centering
\caption{OTTT vs NDOT 在标准 benchmark 上的精度对比}
\begin{tabular}{lcc}
\toprule
\textbf{数据集} & \textbf{OTTT} & \textbf{NDOT} \\
\midrule
CIFAR-10   & 93.73\% & \textbf{94.90\%} \\
CIFAR-100  & 71.11\% & \textbf{76.61\%} \\
DVS-CIFAR10& —       & 77.50\% \\
\bottomrule
\end{tabular}
\label{tab.results}
\end{table}

NDOT 在 CIFAR-100 上相比 OTTT 提升约 +5.5\%，充分证明了用动态神经元动力学替代固定衰减的有效性。

\subsection*{关键结论}

\begin{enumerate}
    \item \textbf{在线训练是可行的}：OTTT 和 NDOT 都证明了，SNN 可以在 O(1) 内存和逐时间步更新的约束下实现有竞争力的精度。
    \item \textbf{物理精确性带来增益}：NDOT 相比 OTTT 的精度提升，完全来自用动态比值 $e[t]$ 替代固定 $\lambda$。这证明了"用物理精确量替换工程近似"这一思路的有效性。
    \item \textbf{在线≠低精度}：NDOT 在 CIFAR-10 上的 94.90\% 已经不输于许多 BPTT 训练的 SNN，证明在线训练可以匹敌离线训练的精度。
\end{enumerate}

\paragraph{讲稿：}
实验结果一目了然：NDOT 在所有 benchmark 上都优于 OTTT，特别是在 CIFAR-100 上有 5.5 个百分点的提升。我想强调的重点是：这个提升是"免费"的——NDOT 的计算量几乎和 OTTT 一样，只是多存了一个 $u[t] - V_{th}s[t]$ 向量，内存仍然是 O(1)。这说明"物理精确性"不是一个纯粹的理论卖点，它在实际精度上有切实的回报。同时也证明了在线训练可以做到和 BPTT 一样好——NDOT 在 CIFAR-10 上的 94.90\% 已经非常接近 BPTT 的 SOTA 水平。

% ============================================================
% Slide 10: 从在线到本地——方法论延伸
% ============================================================
\section{从在线到本地：方法论延伸}

\subsection*{本周总结——我们走到了哪里}

OTTT 和 NDOT 解决了 BPTT 的\textbf{时间维度}瓶颈——实现了前向时间学习和 O(1) 内存。

但如果我们审视梯度公式 $\mathbf{g} \cdot \hat{a}^\top$，注意到空间梯度 $\mathbf{g}_{\mathbf{u}^{l+1}}[t] = \left( \frac{\partial L}{\partial \mathbf{s}^{l+1}[t]} \frac{\partial \mathbf{s}^{l+1}[t]}{\partial \mathbf{u}^{l+1}[t]} \right)^\top$ 仍然需要从输出层逐层反向传播——\textbf{空间维度还没有被"本地化"}。

\subsection*{下一站预览}

下周我们将介绍从"时间本地"迈向"时空全本地"的方法：
\begin{itemize}
    \item \textbf{S-TLLR}：基于 STDP 启发，用三因子学习规则（资格迹 + 学习信号）实现时间本地化，但空间仍依赖 BP/DFA。
    \item \textbf{TESS}：在 S-TLLR 基础上通过 LSG（Local Signal Generation）用固定随机投影在层内生成本地学习信号，彻底切断层间依赖，实现时空全本地化。
\end{itemize}

\paragraph{讲稿：}
最后，让我们回顾一下我们现在的位置。OTTT 和 NDOT 在时间维度上已经实现了突破性的进展——从 O($T$) 到 O(1)，从离线到在线。但我们观察最终的梯度形式 $\mathbf{g} \cdot \hat{a}^\top$，会发现空间梯度 $\mathbf{g}$ 仍然需要从输出层传回来——虽然这个传播不涉及时间，但仍然是"全局的"。

在下周的课程中，我们将探索如何把"本地化"从时间维度扩展到空间维度。S-TLLR 用生物学启发的三因子规则来本地化时间信用分配，TESS 则更进一步，用一个固定随机投影在每一层内部直接生成学习信号，完全不需要层间传播。这将是 SNN 训练方法从"在线"到"全本地"的最后一步跳跃。

\bibliographystyle{IEEEtran}
\bibliography{SNNInformationReconstruction}

\end{document}